% ============================================================================
% Section 2: Background and Related Work scaffold
% ============================================================================
\section{Background and Related Work}
\label{sec:related-work}

\textit{Guidance: Replace these prompts with a focused literature review.
Explain which prior ideas motivate each L\_RAG design choice and identify the
gap addressed by the controlled ablations.}

\subsection{Retrieval-Augmented Generation}
\label{sec:related-rag}

\textit{Guidance: Define the standard query, retriever, context, prompt, LLM,
and answer pipeline. Add a standard RAG diagram and a grounding illustration.}

\subsection{Legal Information Retrieval}
\label{sec:related-legal-ir}

\textit{Guidance: Discuss authority, hierarchy, temporal validity, citation
dependencies, and the cost of retrieval errors in legal settings. Add the
legal-document hierarchy and a legal retrieval example.}

\subsection{Dense, Sparse, and Hybrid Retrieval}
\label{sec:related-hybrid}

\textit{Guidance: Compare dense retrieval, BM25, and rank fusion. Add the
retrieval comparison table and a vocabulary-mismatch example.}

\subsection{Knowledge-Graph-Augmented Retrieval}
\label{sec:related-graph}

\textit{Guidance: Explain how graph structure supplies containment,
neighboring, citation, amendment, and validity context. Add a graph-retrieval
illustration and a vector-versus-graph comparison table.}

\subsection{Reranking and Citation-Grounded Generation}
\label{sec:related-reranking}

\textit{Guidance: Introduce candidate reranking, cross-encoders, citation
grounding, and abstention. Add the two-stage ranking and citation-grounding
illustrations.}

\subsection{Design Requirements for Legal RAG}
\label{sec:related-requirements}

\textit{Guidance: Convert legal reliability requirements into design choices
and evaluation checks. Add a requirement traceability matrix and design
rationale diagram.}
